% Stage 8 restored fragment — 2.8.1 Lorentz reciprocity in operator form (archive lines 10306-10596)
\subsection{Lorentz reciprocity in operator form}
\label{sec:ch281}

On a declared compact-range or anechoic campaign, technicians record frequency, outline, and probe geometry before any harmonic fit is trusted \cite{hansen1988}. Reciprocity links transmission and reception on the same passive structure \cite{stratton1941}. Two-port calibrations depend on it \cite{jackson1999}. Maxwell's retarded export to $\Gamma_R$ remains the governing constraint \cite{maxwell1873}.
Reciprocity is the laboratory fact that passive linear antennas exchange the same complex
transfer when transmit and receive roles swap on a calibrated vector network analyzer
\cite{stratton1941}. The adjoint radiation map extends that symmetry for continuous sources
and traces, enabling source reconstruction and equivalence checks without re-solving unrelated
models \cite{devaney2004}. Reflector feeds, horn coupling, and patch arrays are routinely tuned
by comparing \(S_{21}\) and \(S_{12}\); persistent asymmetry after calibration usually signals a
broken outgoing radiation condition or incorrect impressed-current normalization rather than a
genuine non-reciprocal medium \cite{hansen1988}. Reciprocity cannot be set aside merely because
``only the radiative image matters'': pattern, gain, and scattering matrices on passive hardware
are reciprocity-linked observables tied to the same Maxwell operator
\cite{stratton1941}. Lorentz reciprocity is therefore a consequence of passive
linear Maxwell theory, not an optional post-check applied after the forward solve
\cite{stratton1941}. Two-port calibrations and bistatic radar-cross-section matrices inherit the
same exchange symmetry whenever media and boundaries satisfy the reciprocal class stated below
\cite{jackson1999}. The volume identity that follows states the same swap observed on the
VNA once impressed sources and test fields are paired through the Maxwell adjoint.
\medskip
Section~\ref{sec:ch28} established bilinear exchange symmetry through
Eqs.~\eqref{eq:ch2.section28-reciprocity-chain}--\eqref{eq:ch2.lorentz-reciprocity}.
Subsection~\ref{sec:ch281} supplies the definitive account of Lorentz reciprocity on
\(\mathcal{S}_\omega\) and \(\mathcal{R}_\omega\), including the full proof of dyadic swap
Eq.~\eqref{eq:ch2.dyadic-reciprocity-swap} deferred from Subsection~\ref{sec:ch231}
\cite{stratton1941}. Equation~\eqref{eq:ch2.maxwell-adjoint-pairing}, source dual pairing
Eq.~\eqref{eq:ch2.dual-source-pairing}, outgoing radiation condition
Eqs.~\eqref{eq:ch2.outgoing-radiation-refinement}--\eqref{eq:ch2.outgoing-resolvent-unique},
and outgoing source-to-field integrals of Subsection~\ref{sec:ch251} are presupposed;
Green construction and Sommerfeld uniqueness are not re-opened \cite{stratton1941}. Modal
orthogonality and coupling belong to Subsection~\ref{sec:ch282}; transmit--receive systems to
Subsection~\ref{sec:ch283} \cite{harrington2001}.
Assumptions. Fix \(\omega>0\), Lipschitz domain \(\Omega\) with piecewise smooth interfaces, impressed
sources in \(\mathcal{J}_\Omega\) Eq.~\eqref{eq:ch2.source-space-refined} satisfying
Eq.~\eqref{eq:ch2.source-admissibility}, and outgoing fields in
\(\mathcal{F}_{\mathrm{SF}}\) Eq.~\eqref{eq:ch2.admissible-field-space}
\cite{stratton1941}. Constitutive tensors are linear, local, and passive at
\(\omega\), with symmetric permittivity and permeability in the sense

\begin{equation}
\label{eq:ch2.reciprocal-medium-class}
\overline{\overline{\varepsilon}}(\mathbf{r})=\overline{\overline{\varepsilon}}^{\mathsf{T}}(\mathbf{r}),
\qquad
\overline{\overline{\mu}}(\mathbf{r})=\overline{\overline{\mu}}^{\mathsf{T}}(\mathbf{r}),
\qquad
\overline{\overline{\xi}}=\overline{\overline{\zeta}}=\mathbf{0},
\end{equation}
where \(\overline{\overline{\xi}}\), \(\overline{\overline{\zeta}}\) denote magneto-electric
coupling blocks in the constitutive relation imported from Section~\ref{sec:ch01.1.2}
\cite{stratton1941}. Boundary operator \(\mathcal{B}_\Gamma\) is passive
Eq.~\eqref{eq:ch2.passive-impedance-class} and identical for both excitations
\cite{stratton1941}. Field pairing \(\langle\cdot,\cdot\rangle_{\mathcal{U}}\)
is induced by Eq.~\eqref{eq:ch2.field-energy-norm} on admissible Maxwell pairs
\cite{stratton1941}. regime qualification Eq.~\eqref{eq:ch2.regime-operator-screen} governs transport
interpretation of any bilinear quantity built from \(\Pi_{\mathrm{rad}}\)
\cite{harrington2001}.
Equation~\eqref{eq:ch2.reciprocal-medium-class} is the Chapter~\ref{ch:02} $\mathcal{D}_J=1$
class for \emph{full} Lorentz reciprocity at fixed \(\omega\): symmetric \(\varepsilon\),
\(\mu\) ensure time-reversal symmetry of the harmonic Maxwell system modulo outgoing
selection, while vanishing magneto-electric blocks exclude non-reciprocal gyrotropic or
Tellegen couplings \cite{stratton1941}. Isotropic lossless media are a special
case with scalar \(\varepsilon\), \(\mu\). Lossy symmetric tensors preserve reciprocity
relations stated below with complex conjugation on one slot of each pairing
\cite{stratton1941}. Active or negative-real impedance portions of
\(\mathcal{B}_\Gamma\) break passive symmetry and are excluded from the present theorem
\cite{stratton1941}.
Let \((\E_1,\H_1)=\mathcal{S}_\omega[\Jimp^{(1)}]\) and
\((\E_2,\H_2)=\mathcal{S}_\omega[\Jimp^{(2)}]\) be outgoing admissible solutions for
independent impressed currents \(\Jimp^{(1)},\Jimp^{(2)}\in\mathcal{J}_\Omega\)
\cite{stratton1941}. Here \(\Jimp^{(k)}\) denotes the impressed electric current density in
\(\mathrm{A\,m^{-2}}\), \(\E_k\) and \(\H_k\) are the induced phasor fields in \(\mathrm{V\,m^{-1}}\)
and \(\mathrm{A\,m^{-1}}\), and \(\Gamma\) is the passive boundary on which the surface integral
below is evaluated. Physically, the left member is the bilinear coupling measured when port labels
are swapped on reciprocal hardware; the right member is the surface remainder after integration by
parts. The classical Lorentz reciprocity relation reads
\begin{equation}
\label{eq:ch2.lorentz-volume-reciprocity}
\int_\Omega
\big(\E_1\cdot\Jimp^{(2)*}-\E_2\cdot\Jimp^{(1)*}\big)\,d^{3}r
=
\int_\Gamma
\big(\hat{\mathbf{n}}\times\E_1\cdot\H_2^{*}
-\hat{\mathbf{n}}\times\E_2\cdot\H_1^{*}\big)\,dS,
\end{equation}
when \(\overline{\overline{\varepsilon}}\), \(\overline{\overline{\mu}}\) satisfy
Eq.~\eqref{eq:ch2.reciprocal-medium-class} and \(\mathcal{B}_\Gamma\) is passive
\cite{stratton1941}. Equation~\eqref{eq:ch2.lorentz-volume-reciprocity}
balances impressed-source bilinear forms against a boundary pairing generated by
Eq.~\eqref{eq:ch2.maxwell-adjoint-pairing}: the left member measures how each field couples
to the other's source, while the right member is the surface term that remains after
integration by parts on the Maxwell operator \cite{stratton1941}. Interpretation: when the
boundary term vanishes on admissible outgoing pairs, the port-swap scalar observed on the VNA
equals the volume contraction on the left. On domains where
\(\mathcal{B}_\Gamma\) annihilates the boundary integrand for admissible outgoing pairs---PEC,
homogeneous impedance with Eq.~\eqref{eq:ch2.passive-impedance-class}, and standard
Sommerfeld exterior completions---the surface term vanishes and volume reciprocity closes
\cite{stratton1941}.
Derivation sketch (volume form). Substitute \(\mathbf{u}_\omega^{(1)}=(\E_1,\H_1)\) and \(\mathbf{u}_\omega^{(2)}=(\E_2,\H_2)\)
into Eq.~\eqref{eq:ch2.maxwell-adjoint-pairing} and interchange the two states
\cite{stratton1941}. Symmetry of \(\overline{\overline{\varepsilon}}\),
\(\overline{\overline{\mu}}\) under Eq.~\eqref{eq:ch2.reciprocal-medium-class} makes the
volume pairing sesquilinear form on \(\mathcal{M}_\omega[\mathbf{u}_\omega]\) invariant under
label swap when sources are impressed currents in the Amp\`ere row
\cite{stratton1941}. Subtracting the two adjoint identities eliminates
\(\langle\mathcal{M}_\omega[\mathbf{u}_\omega^{(1)}],\mathbf{u}_\omega^{(2)}\rangle_{\Omega}\)
and yields Eq.~\eqref{eq:ch2.lorentz-volume-reciprocity} with the boundary term displayed
\cite{stratton1941}. Charge continuity Eq.~\eqref{eq:ch1.charge-continuity-harmonic} and
compatibility Eq.~\eqref{eq:ch2.source-compatibility} ensure that impressed \(\Jimp\) alone
specifies the electric drive when \(\rho\) is tied to \(\Jimp\)
\begin{figure}[t]
\centering
\ChOneFigBlock{%
\begin{tikzpicture}[ch1fig, node distance=5mm and 7mm]
 \node[ch1box] (j1) {$\Jimp^{(1)}$};
 \node[ch1box, right=7mm of j1] (s) {$\mathcal{S}_\omega$};
 \node[ch1box, right=7mm of s] (u1) {$\mathbf{u}_1$};
 \node[ch1box, below=5mm of j1] (j2) {$\Jimp^{(2)}$};
 \node[ch1box, below=5mm of s] (u2) {$\mathbf{u}_2$};
 \node[ch1box, below=5mm of u1] (bil) {Bilinear\\pairing};
 \draw[ch1flow] (j1) -- (s) -- (u1);
 \draw[ch1flow] (j2) -- (s) -- (u2);
 \draw[ch1flow, dashed] (u1.south) -- (bil.north);
 \draw[ch1flow, dashed] (j2.east) -- (bil.west);
 \draw[ch1flow, dashed] (u2.east) -- (bil.west);
 \draw[ch1flow, dashed] (j1.east) -- (bil.west);
\end{tikzpicture}%
}
\caption{Subsection~\ref{sec:ch281} Lorentz reciprocity: source labels interchange under
\(\mathcal{S}_\omega\) when media satisfy Eq.~\eqref{eq:ch2.reciprocal-medium-class} and
\(\mathcal{B}_\Gamma\) is passive on the outgoing branch.}
\label{fig:ch2.lorentz-reciprocity-measurement chain}
\end{figure}
Figure interpretation. Sources $\Jimp^{(1)}$, $\Jimp^{(2)}$ induce fields $\mathbf{u}_1$, $\mathbf{u}_2$ through the same outgoing $\mathcal{S}_\omega$. The central bilinear node is symmetric under label interchange when media satisfy Eq.~\eqref{eq:ch2.reciprocal-medium-class} and $\mathcal{B}_\Gamma$ is passive. Dashed return arrows mark that reciprocity is proved on the difference system with identical boundary data, not by comparing unrelated forward and adjoint meshes. Projection to $\mathcal{R}_\omega$ preserves symmetry only after regime qualification Eq.~\eqref{eq:ch2.regime-component-criterion} \cite{stratton1941}. Figure~\ref{fig:ch2.lorentz-reciprocity-measurement chain} is the operator reading of
Eq.~\eqref{eq:ch2.lorentz-volume-reciprocity}: the bilinear functional is symmetric under
label interchange $(\Jimp^{(1)},\mathbf{u}_1)\leftrightarrow(\Jimp^{(2)},\mathbf{u}_2)$ once
boundary terms cancel \cite{stratton1941}.
For vacuum observation points $\mathbf{r}_1$, $\mathbf{r}_2$ with $k_0|\mathbf{r}_1-\mathbf{r}_2|=2.5$, Eq.~\eqref{eq:ch2.dyadic-reciprocity-swap} requires $\overline{\mathbf{G}}_{E,0}(\mathbf{r}_1,\mathbf{r}_2)^{\mathsf{T}}=\overline{\mathbf{G}}_{E,0}(\mathbf{r}_2,\mathbf{r}_1)$ on the outgoing branch. Transverse block entries agree to machine precision; reactive $R^{-2}$, $R^{-3}$ blocks differ by sign under interchange yet cancel in the bilinear contraction when both sides use identical outgoing selectors \cite{stratton1941}. A practical verification proceeds in three steps: confirm the outgoing branch
Eq.~\eqref{eq:ch2.outgoing-green-branch}; test swap symmetry on transverse dyad blocks;
and project to $\mathcal{R}_\omega$ before port calibration \cite{jackson1999}.
\phantomsection\label{thm:ch2.lorentz-reciprocity-operator}
Let \(\mathcal{S}_\omega:\mathcal{J}_\Omega\to\mathcal{F}_\Omega\) be the outgoing solution
map Eq.~\eqref{eq:ch2.maxwell-solution-map} on media satisfying
Eq.~\eqref{eq:ch2.reciprocal-medium-class} with passive \(\mathcal{B}_\Gamma\)
Eq.~\eqref{eq:ch2.passive-impedance-class}. For admissible
\(\Jimp^{(1)},\Jimp^{(2)}\in\mathcal{J}_\Omega\), write
\(\mathbf{u}_k=\mathcal{S}_\omega[\Jimp^{(k)}]\), \(k=1,2\). Then:
\begin{enumerate}
\item[(i)] Volume reciprocity Eq.~\eqref{eq:ch2.lorentz-volume-reciprocity} holds with
vanishing boundary term on admissible outgoing radiation conditions.
\item[(ii)] Operator form: with \(\mathbf{u}_k=(\E_k,\H_k)\) the Maxwell pair induced by
\(\Jimp^{(k)}\) and \(\langle\cdot,\cdot\rangle_{\mathcal{U}}\) the energy inner product on
\(\mathcal{F}_{\mathrm{SF}}\),
\begin{equation}
\label{eq:ch2.lorentz-reciprocity-operator}
\langle \mathbf{u}_1,\,\mathcal{S}_\omega[\Jimp^{(2)}]\rangle_{\mathcal{U}}
=
\langle \mathbf{u}_2,\,\mathcal{S}_\omega[\Jimp^{(1)}]\rangle_{\mathcal{U}}
\end{equation}
when the pairings are defined through impressed-current forcing functionals
Eq.~\eqref{eq:ch2.forcing-functional-operator} and field norm
Eq.~\eqref{eq:ch2.field-energy-norm} on \(\mathcal{F}_{\mathrm{SF}}\).
\item[(iii)] Dyadic realization: outgoing electric dyads satisfy
Eq.~\eqref{eq:ch2.dyadic-reciprocity-swap}.
\item[(iv)] Radiation reciprocity: projecting both sides of (ii) through the self-adjoint
outgoing projector \(\Pi_{\mathrm{rad}}\) yields symmetric coupling on the radiative subspace
\(\mathcal{F}_{\mathrm{rad}}\):
\begin{equation}
\label{eq:ch2.radiation-map-reciprocity}
\langle \mathcal{R}_\omega[\Jimp^{(1)}],\,\mathcal{R}_\omega[\Jimp^{(2)}]\rangle_{\mathcal{U}}
=
\langle \mathcal{R}_\omega[\Jimp^{(2)}],\,\mathcal{R}_\omega[\Jimp^{(1)}]\rangle_{\mathcal{U}}
\end{equation}
on \(\mathcal{F}_{\mathrm{rad}}\) with \(\mathcal{R}_\omega=\Pi_{\mathrm{rad}}\circ\mathcal{S}_\omega\)
Eq.~\eqref{eq:ch2.radiation-operator-def}.
\end{enumerate}
Part (i) follows from Eq.~\eqref{eq:ch2.maxwell-adjoint-pairing} and
Eq.~\eqref{eq:ch2.reciprocal-medium-class} as in the derivation above; outgoing and passive
boundary classes from Subsection~\ref{sec:ch214} cancel the surface integrand
\cite{stratton1941}. Part (ii) identifies the volume impressed-current pairing
with \(\langle\cdot,\cdot\rangle_{\mathcal{U}}\) on \(\mathcal{F}_{\mathrm{SF}}\) through
Eq.~\eqref{eq:ch2.dual-source-pairing} and stability
Eq.~\eqref{eq:ch2.solution-map-stability}, developing Eq.
Eq.~\eqref{eq:ch2.lorentz-reciprocity} \cite{stratton1941}. Part (iii): for fixed
\(\mathbf{r}\), Eq.~\eqref{eq:ch2.electric-dyadic-convolution} and part (i) with impulsive
\(\Jimp^{(k)}\) aligned with Cartesian unit vectors yield column swap
\(\overline{\mathbf{G}}_{E}(\mathbf{r},\mathbf{r}')=
\overline{\mathbf{G}}_{E}^{\mathsf{T}}(\mathbf{r}',\mathbf{r})\) on the outgoing branch
\cite{stratton1941}. Uniqueness exterior radiative uniqueness (Subsection~\ref{sec:ch243})
extends the swap to all admissible \(\mathbf{r},\mathbf{r}'\) \cite{stratton1941}. Part (iv)
applies \(\Pi_{\mathrm{rad}}\) to both sides of (ii); \(\Pi_{\mathrm{rad}}\) is self-adjoint
on \(\mathcal{F}_{\mathrm{rad}}\) under Eq.~\eqref{eq:ch2.field-energy-norm} when outgoing
selection is fixed \cite{stratton1941}. Equation~\eqref{eq:ch2.radiation-map-reciprocity} is the export-projected
form of the VNA symmetry: coupling reported on \(\mathcal{F}_{\mathrm{rad}}\) is unchanged when
port labels swap, while reactive near-zone terms dropped by \(\Pi_{\mathrm{rad}}\) need not be
symmetric even on reciprocal hardware \cite{stratton1941}.
The bilinear functional
\(\mathfrak{B}_\omega(\Jimp^{(1)},\Jimp^{(2)})=
\langle \mathcal{S}_\omega[\Jimp^{(1)}],\,\mathcal{S}_\omega[\Jimp^{(2)}]\rangle_{\mathcal{U}}\)
is symmetric on \(\mathcal{J}_\Omega\times\mathcal{J}_\Omega\) under reciprocal media
\cite{stratton1941}. Algebraically, \(\mathcal{S}_\omega\) is not self-adjoint
on \(\mathcal{J}_\Omega\) with source pairing Eq.~\eqref{eq:ch2.dual-source-pairing}; reciprocity
instead couples the source space to the field space through the \emph{interchanged} labels in
Eq.~\eqref{eq:ch2.lorentz-reciprocity-operator} \cite{stratton1941}. Physically, the far-zone
pattern produced by \(\Jimp^{(1)}\) tested against the field of \(\Jimp^{(2)}\) equals the
swap configuration: a receive pattern computed with \(\Jimp^{(2)}\) as probe matches a
transmit map driven by \(\Jimp^{(1)}\) when hardware and media are reciprocal
\cite{harrington2001}.
Mathematical role of Eq.~\eqref{eq:ch2.radiation-map-reciprocity}. Equation~\eqref{eq:ch2.radiation-map-reciprocity} restricts Eq.~\eqref{eq:ch2.lorentz-reciprocity-operator}
to \(\mathcal{F}_{\mathrm{rad}}\), the subspace on which flux meters and angular patterns are
declared \cite{stratton1941}. Because \(\Pi_{\mathrm{rad}}\) annihilates
non-export components Eq.~\eqref{eq:ch2.rad-null-on-evanescent}, reactive near-zone overlap
does not enter Eq.~\eqref{eq:ch2.radiation-map-reciprocity} even when present in full
Eq.~\eqref{eq:ch2.lorentz-reciprocity-operator} \cite{stratton1941}. regime qualification
Eq.~\eqref{eq:ch2.regime-component-criterion} must still be satisfied before interpreting
Eq.~\eqref{eq:ch2.radiation-map-reciprocity} as a transport channel
\cite{harrington2001}. \phantomsection\label{prop:ch2.dyadic-reciprocity-home}
Under Eq.~\eqref{eq:ch2.reciprocal-medium-class} with outgoing radiation condition, the electric dyad
\(\overline{\mathbf{G}}_{E,\omega}^{\mathrm{(out)}}\) in
Eq.~\eqref{eq:ch2.electric-dyadic-convolution} satisfies
Eq.~\eqref{eq:ch2.dyadic-reciprocity-swap} for all \(\mathbf{r},\mathbf{r}'\) in the
evaluation domain of outgoing source-to-field integrals of Subsection~\ref{sec:ch251}.
Apply Lorentz reciprocity for the solution map (Subsection~\ref{sec:ch281})(iii); this is the proof deferred from
Subsection~\ref{sec:ch231} \cite{stratton1941}. \phantomsection\label{prop:ch2.reciprocity-null-invariance}
If \(\Jimp^{(1)}\) is arbitrary and \(\Jimp^{(2)}=\Jimp^{(1)}+\Jimp^{\mathrm{(NR)}}\) with
\(\Jimp^{\mathrm{(NR)}}\in\ker\mathcal{R}_\omega\), then
Eq.~\eqref{eq:ch2.radiation-map-reciprocity} is unchanged when either argument is replaced by
\(\Jimp^{(2)}\) on the radiative pairing
\(\mathcal{R}_\omega[\Jimp^{(2)}]=\mathcal{R}_\omega[\Jimp^{(1)}]\) by
Eq.~\eqref{eq:ch2.nr-superposition-law} \cite{stratton1941}. Full
Eq.~\eqref{eq:ch2.lorentz-reciprocity-operator} may change by reactive content through
\(\mathcal{S}_\omega[\Jimp^{\mathrm{(NR)}}]\) Eq.~\eqref{eq:ch2.nr-image-of-null-space}
\cite{harrington2001}.
On the outgoing branch, Eq.~\eqref{eq:ch2.lorentz-reciprocity-operator} is consistent with
resolvent symmetry of the truncated Maxwell operator on reciprocal media:
\begin{equation}
\label{eq:ch2.resolvent-reciprocity-swap}
\langle
\mathcal{L}_\omega^{-1}[\mathbf{s}^{(1)}],\,\mathbf{s}^{(2)}
\rangle_{\mathcal{Y}}
=
\langle
\mathcal{L}_\omega^{-1}[\mathbf{s}^{(2)}],\,\mathbf{s}^{(1)}
\rangle_{\mathcal{Y}}
\end{equation}
for compatible forcing pairs \(\mathbf{s}^{(k)}\) when \(\mathcal{L}_\omega\) is symmetric on
\(\mathcal{D}(\mathcal{L}_\omega)\) and \(\mathcal{B}_\Gamma\) is passive
\cite{stratton1941}. Equation~\eqref{eq:ch2.resolvent-reciprocity-swap} is not
re-proved here; it is the operator backbone behind
Eq.~\eqref{eq:ch2.integral-resolvent-equivalence} and Lorentz reciprocity for the solution map (Subsection~\ref{sec:ch281})(ii) \cite{stratton1941}. Numerical MoM matrices
assembled from Eq.~\eqref{eq:ch2.dyadic-reciprocity-swap} inherit swap symmetry only when
quadrature, outgoing branch, and material sampling preserve
Eq.~\eqref{eq:ch2.reciprocal-medium-class} \cite{harrington2001}.
Loss, anisotropy, and broken reciprocity. Complex symmetric \(\overline{\overline{\varepsilon}}\), \(\overline{\overline{\mu}}\) with
\(\Im\{\overline{\overline{\varepsilon}}\}\succeq 0\), \(\Im\{\overline{\overline{\mu}}\}\succeq 0\)
preserve Eqs.~\eqref{eq:ch2.lorentz-volume-reciprocity}--\eqref{eq:ch2.lorentz-reciprocity-operator}
with conjugation on one slot \cite{stratton1941}. Magneto-electric coupling
\(\overline{\overline{\xi}}\neq\mathbf{0}\) breaks Eq.~\eqref{eq:ch2.dyadic-reciprocity-swap}
and requires non-reciprocal Green operators \cite{stratton1941}. Active boundaries
with \(\Re\{Z_s\}<0\) on portions of \(\Gamma_Z\) violate
Eq.~\eqref{eq:ch2.passive-impedance-class} and invalidate Lorentz reciprocity for the solution map (Subsection~\ref{sec:ch281}) \cite{stratton1941}. Incoming rather
than outgoing radiation condition selects a different branch; swap identities hold only when both excitations
use the same branch \cite{stratton1941}. Lorentz reciprocity for the solution map (Subsection~\ref{sec:ch281}) assumes fixed \(\omega\), linear media
Eq.~\eqref{eq:ch2.reciprocal-medium-class}, passive \(\mathcal{B}_\Gamma\), outgoing radiation condition,
and admissible \(\mathcal{J}_\Omega\) \cite{stratton1941}. It does not assert
reciprocity of raw near-zone reactive scans without specifying whether \(\Pi_{\mathrm{rad}}\)
has been applied \cite{stratton1941}. It does not remove \(\ker\mathcal{R}_\omega\)
null-space invariance of radiative reciprocity (Subsection~\ref{sec:ch281}); invisible currents alter full
\(\mathcal{S}_\omega\) pairings while leaving radiative reciprocity unchanged
\cite{harrington2001}. Nonlinear, time-varying, or moving media require the causal chain of
Subsection~\ref{sec:ch223} \cite{stratton1941}.
Subsection~\ref{sec:ch281} established Lorentz reciprocity on \(\mathcal{S}_\omega\) and
\(\mathcal{R}_\omega\) through volume and operator identities
Eqs.~\eqref{eq:ch2.lorentz-volume-reciprocity}--\eqref{eq:ch2.radiation-map-reciprocity}
\cite{stratton1941}. The port-swap symmetry observed on reciprocal hardware is now available as a
constraint on continuous source and field pairings; the next step is to express that constraint on
the modal coordinates fixed in Section~\ref{sec:ch26}.
